A central question in public finance and asset pricing is whether markets efficiently internalize the long-term consequences of infrastructure disinvestment. While a vast literature has established the positive capitalization of new public investment \citep{Oates1969, cellini2010value}, we know surprisingly little about how and when markets price the physical depreciation of the existing public commons. In theory, rational expectations should trigger an immediate adjustment in property values the moment a fiscal decision to reduce maintenance funding is made. In practice, if information frictions or fiscal illusions persist, the economic cost of disinvestment may remain hidden until physical decay becomes tangible. We investigate this disconnect through a quasi-experimental setting in which the underlying fiscal shock --- a local road-maintenance tax cut --- is cleanly identified by close-margin renewal referendums, and in which the underlying public-capital depreciation is directly measured from satellite imagery at sub-county resolution over thirty years.

The empirical setting is well-suited to the question because the deterioration of roads is a {\it visible} phenomenon. Voters and prospective home-buyers can, in principle, observe the consequences of reduced maintenance directly. Whether they do so in a manner that reconciles with the rational-expectations capitalization benchmark is an empirical question that this paper answers in the negative.

\paragraph{Setting and identification.}
We assemble a panel of 3,184 renewal-tax referendums on local road maintenance in Ohio from 1991 through 2021, covering 617 local governments. The timing of each referendum is determined by the predetermined expiration cycle of the prior levy, an average of five years, making the {\it timing} of the vote orthogonal to local economic conditions. The {\it outcome} of close-margin votes is, by construction, unforecastable within a narrow window around the 50-percent threshold. We exploit this discontinuity in a sharp regression-discontinuity design that delivers two empirical objects: the dynamic event-study path of (i) house-price RD coefficients in event time $\tau \in \{-3, \dots, +10\}$ and (ii) road-quality RD coefficients in rolling three-year post-election windows.

The first contribution of the paper is the measurement of $G_t$ --- the underlying physical road-capital stock --- at a resolution and frequency that existing engineering and smartphone-based instruments cannot deliver outside of dense metropolitan settings. We fine-tune a ConvNeXt-V2 vision model on 53,677 labeled satellite road images from \cite{brewer2021} and apply it to the National Agriculture Imagery Program (NAIP) and the Ohio Statewide Imagery Program (OSIP3) image panels. The resulting Road Quality Score, defined on a 0--100 continuous scale, is available year-by-year at sub-county resolution over the full sample window. This is, to our knowledge, the first measurement of dynamic changes in the physical road-capital stock at this resolution over this horizon.

\paragraph{Main empirical findings.}
Three empirical patterns emerge from the analysis. First, jurisdictions that narrowly fail to renew their road-maintenance levy experience an immediate 11-percent decline in maintenance spending --- the direct fiscal margin. Second, road quality declines progressively in the years following the vote, with the rolling-window RD coefficient on the Road Quality Score becoming statistically significant in the $\tau \in [2,4]$ window (a 16-percent decline) and reaching a 23-percent decline at the 10-year horizon. Third, and most importantly, the response of housing prices does not follow the path predicted by rational-expectations capitalization. The pre-period coefficients are economically small and statistically indistinguishable from zero --- which rules out an anticipation story driven by autocorrelated political characteristics. But the price response immediately {\it after} the vote, in $\tau \in [0,3]$, is also indistinguishable from zero. The price gap opens only in $\tau \in [4,9]$, peaks at roughly nine percent of the mean sale price, and partly closes by $\tau = 10$. The total cumulative ten-year discount is approximately \$15,350, or nine percent of the mean median sale price (\$166,000).

\paragraph{The puzzle.}
Under rational expectations, the entire present value of the future amenity decline should have been capitalized at $\tau = 0$, the moment the close-election outcome was realized. The pre-period null effectively rejects an anticipation explanation. The contemporaneous null effectively rejects rational expectations. The price gap opens not when the fiscal shock occurs, but when the {\it physical signal} --- the visible degradation of roads, captured by our ConvNeXt-based measure --- becomes statistically detectable. Approximately one-twentieth of the long-run capitalization arrives in the first three post-election years; nineteen-twentieths arrives only after physical decay is observable.

The synchronization between the visibility of decay and the price response is the central piece of evidence in this paper. It is decisive in distinguishing rational-expectations capitalization from limited-attention and signal-extraction alternatives in which buyers update their belief about the public-capital stock only when an observable signal arrives. We formalize this distinction in Section \ref{sec:fiscal_illusion} and develop a structural counterpart in Section \ref{sec:mech}.

\paragraph{Structural test and welfare.}
We embed the empirical wedge in a dynamic general-equilibrium model with a Rioja-style road-depreciation channel. Under perfect foresight --- the rational-expectations benchmark --- the model overshoots the data in the short run, generating an immediate price drop that is decisively rejected by the empirical event study. We then introduce a Bayesian signal-extraction microfoundation in which households observe the public-capital stock with noise that shrinks endogenously with the magnitude of physical decay. The microfoundation nests the rational-expectations benchmark as a knife-edge restriction ($\bar{\sigma} = 0$) and identifies the friction parameter $\bar{\sigma}$ from a single transparent moment in the data: the ratio of the cumulative discounted price gap accumulated through $\tau = 3$ to the long-run capitalization at $\tau \geq 9$. The calibrated belief-updating model lies within the empirical 95-percent confidence band throughout the event-time window and reproduces the cross-sectional heterogeneity by urban density and by housing-price quantile without additional free parameters.

The implied welfare loss is substantial. Capitalizing the average per-household tax savings (\$76 per year) as a perpetuity at the user-cost rate of 0.04 yields a present value of approximately \$1,900. The corresponding capitalized wealth loss is \$15,350. The ratio --- approximately eight to one --- is the {\it capitalization gap} we report in the body of the paper. Translated into the marginal-value-of-public-funds (MVPF) framework of \citet{hendren2020unified}, the policy delivers an MVPF of $-1.4$: every dollar of revenue forgone by the government generates approximately \$1.40 of net welfare loss to households. Under the belief-updating microfoundation, this number is best interpreted as a conservative {\it choice-based} measure; the {\it experience-based} counterpart, evaluated under perfect foresight, lies in the range $[-2.1,\, -1.6]$.

\paragraph{Contributions.}
The paper makes three contributions. First, it provides the first credibly-identified estimates of the capitalization of public-capital {\it depreciation} into property values --- a literature in which positive shocks (new bonds, new highways) have dominated and the negative-shock case has been substantially understudied despite its first-order policy importance for an aging U.S. infrastructure stock \citep{asce2025report, ramey2021infrastructure}. Second, the paper introduces a novel computer-vision-based measurement of physical road quality at sub-county resolution over thirty years, complementing the smartphone-acceleration approach of \cite{currier2023} and the engineering-survey benchmarks (IRI, PCI) which are sparse outside large metropolitan areas. Third, the paper develops a formal test and rejection of rational-expectations capitalization in a salient local asset market and provides a tractable structural alternative --- a belief-updating microfoundation --- that quantitatively rationalizes the empirical pattern and delivers welfare-relevant policy implications, including pre-vote disclosure of road-condition assessments \citep[in the spirit of][]{brasington2024fire, chetty2009salience}.

\paragraph{Roadmap.}
The remainder of the paper is organized as follows. Section \ref{sec:lit_review} positions the contribution relative to the literatures on public-capital capitalization, information frictions in housing markets, and computer-vision-based economic measurement. Section \ref{sec:data} describes the data and the fine-tuned vision model. Section \ref{sec:empirical} develops the close-election regression-discontinuity design and formalizes the anticipation test. Section \ref{sec:result} presents the empirical results. Section \ref{sec:fiscal_illusion} formalizes the rational-expectations benchmark, documents the empirical wedge, and identifies the candidate mechanisms. Section \ref{sec:mech} develops the structural counterpart, including the belief-updating microfoundation. Section \ref{sec:mvpf} reports the welfare analysis. Section \ref{sec:conclusion} concludes. Technical derivations, additional robustness checks, and the vision-model fine-tuning details are deferred to the online appendix.
